\subsubsection{(Simple) Sort Merge Join (SMJ)}
We first sort $R$ and $S$ on the join attributes with external merge sort (see Section~\ref{sec:external-sort}). We then read the sorted files and merge.
This is of course not great for performance, but if the relations are already sorted, we can skip the first step.

The cost (if there are \bi{no duplicates}) is $\texttt{Sort}(S) + \texttt{Sort}(R) + P_R + P_S$ I/Os.

If there are duplicates (i.e. the keys are not unique, say we are joining on a column that doesn't have the UNIQUE constraint)
then we need to move the pointer for the join back on the second relation to make sure we don't miss any tuples.

The cost now of course isn't linear anymore and worst case quadratic, $\tco{\texttt{Sort}(S) + \texttt{Sort}(R) + P_R \cdot P_S}$.

Compared to BNLJ, its performance doesn't depend on the amount of memory available, but tends to be a tad slower than BNLJ, if BNLJ is provided with ample memory.
